\section{Qu'est-ce que la recherche d'information (Information Retrieval)\,?}

\begin{frame}{\textit{Information Retrieval}, définition}

\begin{block}{Une définition}
La \alert{Recherche d'Information} (\textit{Information Retrieval}, IR) consiste à
trouver des \alert{documents} peu ou faiblement structurés, dans une grande \alert{collection},
en fonction d'un \alert{besoin d'information}.
\end{block}

Prend peu à peu la prédominance sur les techniques de bases de données.

\begin{itemize} 
  \item Recherche sur le Web. Utilisée quotidiennement par des milliards d'utilisateurs.
   \item Recherche dans votre boîte mail.
   \item Recherche sur votre ordinateur (\textit{Spotlight}).
    \item Recherche dans une base documentaire, publique ou privée.
\end{itemize}

À bien distinguer d'un recherche type ``base de données''\,: requête structurée, données structurées,
réponse \guill{exacte}.
\end{frame}

\section{Evaluation d'un moteur de recherche}

\begin{frame}{Spécificités}

Contrairement à une base de données (SQL), le résultat dépend de \alert{l'interprétation}
d'un besoin.

\vfill

On ne peut jamais dire qu'un résultat est totalement exact (ou totalement faux).

\vfill

Deux notions importantes:

\begin{itemize}
  \item \alert{Faux positifs}: ce sont les documents \alert{non pertinents} inclus dans le résultat;
    ils ont été sélectionnés à tort.
  \item \alert{Faux négatifs}: ce sont les documents \alert{pertinents} qui \alert{ne sont pas}
    inclus dans le résultat. 
\end{itemize}

\vfill

En général, chercher à réduire les faux positifs entraîne l'augmentation des faux négatifs, 
et réciproquement. 

\end{frame}

\begin{frame}{Evaluation de la pertinence: précision et rappel}

La \alert{précision} mesure la fraction des vrais positifs dans le résultat $r$.

\medskip

  Si on note $P_r$ et $FP_r$ le nombre de vrais
   et de faux positifs dans $r$, alors
   
  $$\text{précision} = \frac{P_r}{P_r + FP_r} = \frac{P_r}{|r|}$$

\vfill

   Le \alert{rappel} mesure la fraction de faux négatifs. 
   
\medskip
   Si on note $VP$ le
   nombre de vrais positifs \alert{dans la collection} (supposé connu), alors  le rappel est:
   
   $$\frac{P_r}{VP}$$
   
   \vfill
   
   L'évaluation d'un système de RI est difficile: implique des tests rigoureux avec des utilisateurs
   sur un échantillon.

\end{frame}

\section{La recherche plein texte}


\begin{frame}{Notre exemple de base}

La RI permet d'effectuer des recherches \alert{booléennes} et des
recherches \alert{plein texte}. Dans ce dernier cas, il faut \alert{analyser}
des documents textuels.

\vspace{.5cm}

Un ensemble (modeste) nous servira de guide.

\vspace{.5cm}


\begin{tabular}{>{\color{structure.fg}}ll}
     $d_1$&Le loup est dans la bergerie.\\
     $d_2$&Le loup et le trois petits cochons.\\
     $d_3$&Les moutons sont dans la bergerie.\\
     $d_4$&Spider Cochon, Spider Cochon, il peut marcher au plafond.\\
     $d_5$&Un loup a mangé un mouton, les autres loups sont restés dans la bergerie.\\
     $d_6$&Il y a trois moutons dans le pré, et un mouton dans la gueule du loup.\\
     $d_7$&Le cochon est à 12 le Kg, le mouton à  10 E/Kg.\\
     $d_8$&Les trois petits loups et le grand méchant cochon.\\
\end{tabular}


\end{frame}

\begin{frame}{Le besoin et la solution}

On veut chercher tous les documents parlant de loups, de moutons mais pas de bergerie (c’est le besoin).


\medskip

\alert{Parcourir tous les documents?}
\begin{itemize}
  \item  potentiellement long;
  \item critère "pas de bergerie" n'est pas facile à traiter;
  \item  autres types de recherche ("le mot 'loup' doit être \alert{près} du mot 'mouton'") sont
    difficiles;
 \item comment \alert{classer} par pertinence les documents trouvés?
\end{itemize}

\medskip

Structure spécialisée: la \alert{matrice d'incidence} et surtout son inversion.

\end{frame}

\begin{frame}{Matrice avec documents en ligne}

On sélectionne un ensemble de mots (ou \alert{termes}), constituant notre
\alert{vocabulaire} (ou \alert{dictionnaire}).

\vfill

Documents en ligne, termes en colonnes. Dans chaque cellule: 1 si le terme est dans
le document, 0 sinon.

\medskip

\figSlideWithSize{matrice-incidence1}{10cm}

\end{frame}


\begin{frame}{Pour effectuer la recherche}

On prend les vecteurs binaires des \alert{termes} (les colonnes).

\begin{itemize}
  \item Loup:      11001101
  \item Mouton:    00101110
  \item Bergerie:    01010011
\end{itemize}

\medskip

Puis:

\begin{itemize}
  \item ET logique sur les vecteurs de Loup et Mouton, on obtient  00001100.
  \item ET logique  avec le \alert{complément} du vecteur de Bergerie (01010111)
\end{itemize}

On obtient 00000100, d'où on déduit que la réponse est limitée au document $d_6$.

\medskip

Opération binaires \alert{très efficace}, mais...

\end{frame}


\begin{frame}{Passons à grande échelle}

Quelques hypothèses:
\begin{itemize}
 \item Un million de documents, mille mots chacun en moyenne.
 \item Disons 6 octets par mot, soit 6 GO (pas très gros en fait!)
 \item Disons 500\,000 termes \alert{distincts}
\end{itemize}

$\Rightarrow$ la matrice a $500 \times 10^9$ bits, soit 62 GO.

\vfill

Ne tient pas en mémoire, ce qui va beaucoup compliquer les choses....

\vfill

\alert{Comment faire mieux\,?}

\end{frame}

\begin{frame}{On peut faire mieux}

  \begin{minipage}[c]{.46\linewidth}
       Il vaut mieux avoir les termes en ligne pour disposer des vecteurs dans une zone
mémoire contigue.

\vspace*{1cm}

On parle de matrice inversée, et de \alert{liste inversée}.
Une liste par terme; dans chaque liste, 1 pour les documents contenant le terme.
   \end{minipage} \hfill
   \begin{minipage}[c]{.46\linewidth}
\figSlideWithSize{inverted-index-bool}{4.5cm}
   \end{minipage}
   

\begin{block}{Important}
Les mises à jour beaucoup plus difficiles car elles impliquent la réorganisation
d'une liste compacte. Prix à payer pour l'efficacité en \alert{lecture} (recherche).
\end{block}
\end{frame}


\begin{frame}{On peut encore faire mieux}

La matrice est \alert{creuse}: il n'y a que $10^9$ positions avec des 1, soit un sur 500.

\medskip


\figSlideWithSize{inverted-index}{6cm}

\medskip

\begin{block}{Essentiel}
On place dans les cellules \alert{l'identifiant} du document. De plus chaque liste est
\alert{triée} sur l'identifiant du document.
\end{block}

\end{frame}



\begin{frame}{Index inversé}

La structure utilisée dans \alert{tous}  les moteurs de recherche.

\begin{itemize}
 \item Un \alert{répertoire} contient tous les \alert{termes}.
 \item Une \alert{liste} (inversée) est associée à chaque \alert{terme}, \alert{triée} par docId.
 \item  Chaque élément de la liste est appelé une \alert{entrée}.
\end{itemize}

\vfill

\alert{Concept}\,: la notion de \alert{terme} (\alert{token} en anglais) est différente de celle de ``mot''. À revoir.
\vfill

\alert{Vocabulaire}\,: le répertoire est parfois appelé \emph{dictionnaire}\,; les listes sont des \emph{posting list} en anglais;
les entrées sont des \emph{postings}.

\vfill

\alert{Efficacité}\,: le répertoire devrait toujours être en mémoire\,; les listes, autant que possible en mémoire, sinon
fichiers contigus sur le disque.
\end{frame}

\section{Opération de recherche}


\begin{frame}{Traitement d'une recherche par \emph{fusion} de listes triées}

Recherchons les documents parlant de loup et de mouton.

\vfill

Algorithme par \alert{fusion}: on parcourt \alert{séquentiellement} les deux listes.

\smallskip

On compare à chaque étape les docId\,; s'ils sont égaux\,: \alert{trouvé}!

\smallskip

On avance sur la liste du plus petit docId.

\only<1>{
\figSlideWithSize{brutus-et-cesar}{7cm}
}
\only<2>{
\figSlideWithSize{brutus-et-cesar-2}{7cm}
}
\only<3>{
\figSlideWithSize{brutus-et-cesar-3}{7cm}
}
\only<4>{
\figSlideWithSize{brutus-et-cesar-4}{7cm}
}
\only<5>{
\figSlideWithSize{brutus-et-cesar-5}{7cm}
}
\only<6>{
\figSlideWithSize{brutus-et-cesar-6}{7cm}
}
\only<7>{
\figSlideWithSize{brutus-et-cesar-7}{7cm}
}
\only<8>{
\figSlideWithSize{brutus-et-cesar-8}{7cm}
}
\only<9>{
\figSlideWithSize{brutus-et-cesar-9}{7cm}
}

\pause
\vfill

\alert{Un seul parcours suffit}\,: recherche linéaire, parcours séquentiel. Pas mieux.

\vfill

 C'est une recherche dite \alert{Booléenne}\,: pas de classement, résultat exact.
 
 %\vfill
 
 %Question\,: et si je ne veux pas Calpurnia\,? Comment je fais\,?
 
\end{frame}


\begin{frame}{Premier bilan}

En résumé: index inversé, tri et compaction, parcours linéaire très rapide sont
les fondements techniques de la RI.

\begin{itemize}
 \item Permet d'effectuer des recherches \alert{puissantes} (combinaison de critères)
     et \alert{flexibles}.
 \item \alert{Garantissent une très grande efficacité.}
\end{itemize}
\medskip

Que reste-t-il à faire\,?

\begin{description}
 \item[Quels termes\,?] Quels termes indexe-t-on\,? Beaucoup moins facile que ça n'en a l'air...
 \item[Quelles requêtes\,?] De la plus simple (``sac de mots'') à plus 
                 structurée (index struturé, requêtes Boolénnes, recherche de phrases).
   \item[Performance.] Construction, compression, distribution, optimisation des accès,  mises à jour, etc.
  \item[Classement\,?] Si j'ai des millions de documents dans le résultat, je veux les classer. Comment\,?
\end{description}


\end{frame}
